% Lösungen Einheit 4 von 4 – Mit Größen rechnen im Sachzusammenhang (zusammenbau v0.1)
\einheitenkopf{Einheit 4 von 4 · Mit Größen rechnen im Sachzusammenhang}

%% TODO Lösung der Erklärzeile 34g) fehlt in der Bank
\erg{34}{a) $1{,}2$ m \quad b) $1{,}4 + 0{,}65 = 2{,}05$ km \quad c) $2 + 0{,}75 = 2{,}75$ kg \quad d) $3{,}2 + 0{,}45 = 3{,}65$ m \quad e) $4{,}20 + 0{,}85 = 5{,}05$ € \quad f) $1{,}5 + 0{,}33 = 1{,}83$ l \quad g) \ldots \quad h) $8$ cm \quad i) $30 \cdot 2{,}7 = 81$ g \quad j) $445 : 8{,}9 = 50$ cm³ \quad k) $30\,\%$ von $2$ m³ $= 0{,}6$ m³ $= 600$ kg; $600 + 85 = 685$ kg \quad l) $18 : 6 = 3$ l \quad m) $2 \cdot 4 \cdot 2{,}6 = 20{,}8$ m², zweimal: $41{,}6$ m²; $41{,}6 : 8 = 5{,}2$ l \quad n) die $2$-l-Dose ($1{,}3 > 1$) \quad o) Die Kilometer wurden nicht umgerechnet: $1{,}6$ km $= 1\,600$ m; $1\,600 + 750 = 2\,350$ m \quad p) $V = 3\,500 : 8{,}4 \approx 416{,}7$ cm³ \quad q) Fläche $(2 \cdot 0{,}18 + 4 \cdot 0{,}12) \cdot 2 = 1{,}68$ m²; Bedarf $1{,}68 : 8 = 0{,}21$ l $= 210$ ml; $125$ ml reichen nicht, $250$ ml reichen – Max hat recht \quad r) Fehler bei der Umrechnung: $1$ € $= 100$ ct, also $744$ € $= 74\,400$ ct, nicht $744\,000$ ct; $74\,400 : 120\,000 = 0{,}62$ ct – jeder Kilometer wird nur etwa $0{,}6$ ct teurer}
%% TODO Lösung der Erklärzeile 35g) fehlt in der Bank
\erg{35}{a) nein \quad b) $23 + 17 = 40$ Stück, $40 \cdot 2 = 80$ € \quad c) $48 + 36 = 84$ Karten, $84 \cdot 7 = 588$ € \quad d) $120 + 95 + 105 = 320$ Brötchen, $320 \cdot 0{,}45 = 144$ € \quad e) $18 + 25 + 21 = 64$ Stück, $64 \cdot 12 = 768$ € \quad f) $136 + 89 = 225$ Karten, $225 \cdot 15{,}50 = 3\,487{,}50$ € \quad g) \ldots \quad h) $(22 + 17) \cdot 6 = 234$ €; $(15 + 9 + 12) \cdot 4 = 144$ €; zusammen $= 378$ € \quad i) $7\,500$ cm² $= 0{,}75$ m²; $0{,}75 \cdot 30 = 22{,}50$ € \quad j) $5 \cdot 25 + 3 \cdot 25 = 200$ cm; Verschnitt $240 - 200 = 40$ cm; $40 : 240 \approx 0{,}167 = 16{,}7\,\%$}
\erg{36}{a) $540 : 200 = 2{,}7$ g/cm³}
\erg{37}{a) $3 \cdot 0{,}85 = 2{,}55$ €; $2{,}55 + 2{,}40 = 4{,}95$ €}
\erg{38}{a) Tim hat mit tausend statt mit hundert umgerechnet; $1$ € $= 100$ ct: $7{,}25 \cdot 100 = 725$, also $7{,}25$ € $= 725$ ct \quad b) Lina hat die Preise gemittelt, obwohl verschieden viele Getränke verkauft wurden; richtig: $45 \cdot 3 + 25 \cdot 1 = 135 + 25 = 160$ € \quad c) In Metern sind es $7$ m – eine kleine ganze Zahl, die man sich vorstellen kann; $0{,}007$ km hat zwei Nullen hinter dem Komma. \quad d) Der Durchschnittspreis zählt beide Preise gleich oft; es wurden aber viel mehr teure Karten verkauft – jede Preisstufe muss mit ihrer eigenen Stückzahl malgenommen werden. \quad e) $30 \cdot 30 = 900$ cm² $= 0{,}09$ m²; $6{,}3 : 0{,}09 = 70$ Fliesen}
